% ===========================================================================
\section{Geometry and isoperimetry}\label{sec:geometry}
\warmth{0}\quad\whisper{Where ``concentration of measure'' gets its name: on a metric measure space, a Lipschitz function is almost constant.}

\begin{strand}
In one line: a viewpoint with no MGF at all. On a metric measure space, every $1$-Lipschitz function
\keyword{hugs its median}; this is equivalent to an \keyword{isoperimetric} phenomenon, a set of measure
$\ge\tfrac12$, blown up by $t$, swallows $1-e^{-ct^2}$ of the measure. High-dimensional spheres and Gaussian
space behave exactly so; this is the geometric root of ``a random vector's length is almost determined'' and
``a random projection preserves distance''.
\end{strand}

\block{The paradigm}
\trigger{$f$ is an $L$-Lipschitz function on a space with \keyword{concentration function} $\alpha(t)$.}
Then $\Pr[|f-\operatorname{Med}f|>t]\le 2\,\alpha(t/L)$. Three pillar spaces:
\[
  \begin{aligned}
    \text{sphere } S^{n-1}:&\ \alpha(t)\le e^{-(n-1)t^2/2};\\
    \text{Gaussian } \gamma_n:&\ \alpha(t)\le e^{-t^2/2}\quad(\text{\keyword{dimension-free}});\\
    \text{Hamming cube}:&\ \alpha(t)\le e^{-2t^2/n}.
  \end{aligned}
\]

\begin{strand}
Three-line heuristic: (1) concentration $\equiv$ isoperimetry: the tail bound is the complement of ``the
$t$-blow-up of a half-measure set''; (2) median and mean are interchangeable (differ by $O(L)$, same order);
(3) the Gaussian case being \emph{dimension-free} is the key miracle, it makes infinite-/high-dimensional
processes (§\ref{sec:chaining}) possible.
\end{strand}

\block{Cheat-sheet: geometric concentration}
\noindent\begin{tabularx}{\linewidth}{@{}l l L@{}}
\toprule
{\color{indigo}Theorem} & {\color{indigo}Space} & \multicolumn{1}{l}{\color{indigo}Statement / when to use}\\
\midrule
Lévy's lemma    & sphere $S^{n-1}$          & $1$-Lip $f$: $\Pr[|f-\operatorname{Med}f|>t]\le 2e^{-(n-1)t^2/2}$\\
Gaussian conc.\ & $\gamma_n=N(0,I_n)$       & $L$-Lip $f$: $\Pr[|f-\E f|>t]\le 2e^{-t^2/2L^2}$, \emph{dimension-free}\\
Borell–TIS      & Gaussian process         & $\sup_t X_t$ concentrates around its supremum, variance $=\sup_t\Var X_t$\\
Gaussian isoper.& $\gamma_n$                & half-spaces are extremal; exact blow-up in $\Phi^{-1}$ form\\
\bottomrule
\end{tabularx}

\block{Main results \& proof skeletons}

\begin{theorem}[Lévy's lemma / spherical concentration]\label{thm:levy}
$f:S^{n-1}\to\R$ is $1$-Lipschitz. Then $\Pr[|f-\operatorname{Med}f|>t]\le 2e^{-(n-1)t^2/2}$.
\end{theorem}
\begin{proof}
Skeleton: spherical \keyword{isoperimetry} (Lévy–Schmidt: geodesic balls are extremal) $\Rightarrow$ for a
half-measure set $A$, the $t$-neighborhood satisfies $\sigma(A_t)\ge 1-e^{-(n-1)t^2/2}$; take $A=\{f\le\operatorname{Med}f\}$,
and Lipschitz pushes $\{f>\operatorname{Med}f+t\}$ outside $A_t$. \TODO{fill the standard isoperimetry $\Rightarrow$ concentration translation; cite spherical isoperimetry (do not reprove)}
\end{proof}

\begin{theorem}[Gaussian concentration / Borell–TIS]\label{thm:gaussian-conc}
$f:\R^n\to\R$ is $L$-Lipschitz, $X\sim N(0,I_n)$. Then $\Pr[|f(X)-\E f|>t]\le 2e^{-t^2/2L^2}$, with constants independent of $n$.
\end{theorem}
\begin{proof}
Either path: \emph{(a)} Gaussian log-Sobolev (Thm~\ref{thm:gaussian-lsi}) $+$ Herbst (Thm~\ref{thm:herbst}), with smooth
approximation of $f$; or \emph{(b)} Gaussian isoperimetry (half-space extremality) gives the blow-up directly.
\TODO{take (a): fill the Lipschitz-gradient equivalence $\|\nabla f\|\le L$ a.e.\ + smoothing}
\end{proof}

\begin{example}[A random vector's length is almost determined]\label{ex:norm-conc}
$x\mapsto\|x\|$ is $1$-Lipschitz $\Rightarrow$ for $X\sim N(0,I_n)$, $\|X\|$ concentrates near $\sqrt n$ with
fluctuation $O(1)$ (independent of $n$). This is the one-line proof of the ``thin-shell'' phenomenon.
\TODO{fill $\E\|X\|=\sqrt n(1-o(1))$ and substitute}
\end{example}

\begin{strand}
Why it goes this way: the MGF method (§\ref{sec:chernoff}) answers ``sums''; the geometric method answers
``\emph{any} Lipschitz function'', dimension-free in the Gaussian case, at the cost of needing a ``hard'' input
like isoperimetry / log-Sobolev. The two paths \emph{meet} at Gaussian concentration: the LSI of §\ref{sec:entropy}
is precisely the algebraic proof of Borell–TIS.
\end{strand}

\loose{Median vs mean: they differ by $O(L)$, but exact isoperimetry is most natural with the median. When must they be distinguished? (heavy tails, asymmetric $f$.)}%
\warp{borelltis} Borell–TIS (concentration of the supremum of a Gaussian process) is the ``concentration half'' of the chaining in §\ref{sec:chaining}, chaining controls the \emph{expectation}, this theorem the \emph{fluctuation}. \pick{borelltis}
